\documentclass[11pt,a4paper]{article}

% ---------------------------------------------------------
% Modern RDCC Preamble (v48.0)
% ---------------------------------------------------------
\usepackage[utf8]{inputenc}
\usepackage[T1]{fontenc}
\usepackage{lmodern}
\usepackage{amsmath, amssymb, amsfonts}
\usepackage{physics}
\usepackage{graphicx}
\usepackage{float}
\usepackage{cite}
\usepackage{hyperref}
\usepackage{geometry}
\geometry{margin=2.4cm}

\hypersetup{
    colorlinks=true,
    linkcolor=blue,
    citecolor=blue,
    urlcolor=blue
}

\title{\textbf{Companion II}\\[4pt]
\large Tensor Interference and Log-Periodic Gravitational Waves\\
in Relaxation-Driven Cyclic Cosmology (RDCC)}
\author{Michael Lehmann\\
\small Independent Researcher}
\date{June 2026}

\begin{document}

\maketitle

\begin{abstract}
Relaxation-Driven Cyclic Cosmology (RDCC) predicts a distinctive tensor signature
arising from the CPT-symmetric structure of the Universe. In this framework, tensor
perturbations do not originate from stochastic amplification at horizon exit, but from
CPT-consistent matching of tensor modes across a nonsingular bounce. The bipartite
Hilbert-space structure implies that visible and mirror tensor sectors interfere, producing
a log-periodic modulation of the stochastic gravitational-wave background (SGWB).
This Companion develops the structural origin of this modulation, its dependence on
the infrared relaxation parameter $\alpha_{\mathrm{IR}}$, and its observational consequences.
\end{abstract}

\section*{Executive Summary}

Tensor perturbations in RDCC arise from the interplay between three structural elements:
(i) the CPT-symmetric global state, (ii) the bipartite Hilbert-space factorization into visible
and mirror sectors, and (iii) the relaxation dynamics governed by the infrared parameter
$\alpha_{\mathrm{IR}}$. Together, these ingredients imply that tensor modes in the two sectors
acquire a small logarithmic phase difference across the bounce. Their interference produces
a universal log-periodic modulation of the SGWB.

The modulation frequency is proportional to the infrared relaxation scale,


\[
    \omega_{\mathrm{mod}} \propto \alpha_{\mathrm{IR}},
\]


reflecting the one-parameter nature of RDCC. The resulting tensor spectrum exhibits a
blue tilt, a millihertz-scale peak, and a discrete scale invariance encoded in oscillations
in $\ln k$ or $\ln f$. These features provide a sharply falsifiable prediction for upcoming
gravitational-wave observatories such as LISA, DECIGO, BBO, and the Einstein Telescope.

This Companion presents the modernized v48.0 formulation of the tensor sector, including
the global tensor equation, the bounce matching conditions, the CPT-induced phase
structure, and the derivation of the log-periodic modulation. Only two figures are used:
a schematic of tensor-sector interference and a one-dimensional illustration of the
log-periodic modulation.

\section{Introduction}

Relaxation-Driven Cyclic Cosmology (RDCC) describes the Universe as a globally
CPT-symmetric quantum state with a bipartite Hilbert-space structure. Gravity couples
to the sum of the visible and mirror stress-energy tensors, while all gauge interactions
remain strictly sectoral. This structure implies that tensor perturbations propagate on a
shared metric but originate from two conjugate sectors whose dynamics are correlated
through the relaxation parameter $\alpha(a)$.

At the cosmological bounce, the global state is maximally coherent and satisfies a
CPT-symmetric matching condition. Tensor modes in the visible and mirror sectors are
therefore related by time reversal up to a small relaxation-induced phase shift. As the
Universe evolves away from the bounce, the relaxation parameter grows toward its
infrared fixed point $\alpha_{\mathrm{IR}}$, inducing a logarithmic phase difference between the
two tensor branches.

The interference between these CPT-conjugate tensor modes produces a universal
log-periodic modulation of the stochastic gravitational-wave background (SGWB). This
modulation is not a phenomenological addition but a structural consequence of:
\begin{enumerate}
    \item the bipartite Hilbert-space factorization,
    \item the CPT-symmetric bounce matching,
    \item and the relaxation dynamics governed by $\alpha_{\mathrm{IR}}$.
\end{enumerate}

The modulation frequency is proportional to the infrared relaxation scale,


\[
    \omega_{\mathrm{mod}} \propto \alpha_{\mathrm{IR}},
\]


reflecting the one-parameter nature of RDCC. The resulting tensor spectrum exhibits
three characteristic features:
\begin{itemize}
    \item a blue tilt originating from the cuscuton-induced bounce,
    \item a millihertz-scale peak set by the bounce scale,
    \item and a discrete scale invariance encoded in oscillations in $\ln k$ or $\ln f$.
\end{itemize}

These signatures span more than twenty decades in frequency and provide a sharply
falsifiable prediction for upcoming gravitational-wave observatories. The purpose of this
Companion is to present the modern v48.0 formulation of the tensor sector, including
the global tensor equation, the bounce matching conditions, the CPT-induced phase
structure, and the derivation of the log-periodic modulation. Only two figures are used:
a schematic of tensor-sector interference and a one-dimensional illustration of the
log-periodic modulation.

\section{Tensor Equations in a CPT-Symmetric Two-Sector Universe}

The tensor sector of RDCC inherits its structure from the bipartite Hilbert-space
factorization and the global Einstein equations. Gravity couples universally to the sum
of the visible and mirror stress-energy tensors, while gauge interactions remain strictly
sectoral. Tensor perturbations therefore propagate on a single shared metric but receive
contributions from two CPT-conjugate sectors.

We decompose the metric as
\begin{equation}
    g_{\mu\nu} = \bar{g}_{\mu\nu} + h_{\mu\nu},
\end{equation}
where $\bar{g}_{\mu\nu}$ is the spatially flat FLRW background and $h_{\mu\nu}$ is a
transverse-traceless tensor perturbation satisfying
\begin{equation}
    \nabla^{i} h_{ij} = 0,
    \qquad
    h^{i}{}_{i} = 0.
\end{equation}

Linearizing the global Einstein equations yields
\begin{equation}
    \delta G_{ij}
    = 8\pi G \left( \delta T^{(+)}_{ij} + \delta T^{(-)}_{ij} \right),
\end{equation}
where the superscripts $(+)$ and $(-)$ denote the visible and mirror sectors,
respectively. For tensor modes, the only relevant contribution is the anisotropic stress,
\begin{equation}
    \delta T^{(\pm)}_{ij} = \pi^{(\pm)}_{ij}.
\end{equation}

The global tensor equation becomes
\begin{equation}
    \ddot{h}_{ij} + 3H \dot{h}_{ij} - \frac{\nabla^{2}}{a^{2}} h_{ij}
    = 16\pi G \left( \pi^{(+)}_{ij} + \pi^{(-)}_{ij} \right).
\end{equation}

A sectoral observer perceives only the projection onto their own branch. For the visible
sector, this yields
\begin{equation}
    \ddot{h}^{(+)}_{ij}
    + 3H \dot{h}^{(+)}_{ij}
    - \frac{\nabla^{2}}{a^{2}} h^{(+)}_{ij}
    = 16\pi G_{\mathrm{eff}}\, \pi^{(+)}_{ij},
\end{equation}
with an effective Newton constant
\begin{equation}
    G_{\mathrm{eff}}
    = G \left( 1 + \frac{\rho_{-}}{\rho_{+}} \right),
\end{equation}
where $\rho_{+}$ and $\rho_{-}$ denote the energy densities of the visible and mirror
sectors. This expression reflects the fact that gravity couples to the sum of the two
stress-energy tensors, while a sectoral observer perceives only their own anisotropic
stress.

The difference between global and sectoral propagation is the structural origin of the
interference pattern that leads to log-periodic modulation. The global tensor mode
contains contributions from both sectors, while the sectoral mode is a projection onto a
single branch. The relaxation dynamics encoded in $\alpha(a)$ induce a small logarithmic
phase difference between the two tensor branches, which becomes observationally
relevant once the Universe evolves away from the bounce.

In the next section we analyze how the CPT-symmetric bounce matching conditions
determine the relative phase between the visible and mirror tensor modes and how this
phase structure leads to discrete scale invariance in the tensor spectrum.

\section{Bounce Matching and CPT Conditions}

The RDCC bounce is a nonsingular transition between contraction and expansion,
generated by a cuscuton-induced modification of the Friedmann equation. At the
bounce time $t_{\mathrm{b}}$, the background satisfies
\begin{equation}
    H(t_{\mathrm{b}}) = 0,
    \qquad
    \alpha(t_{\mathrm{b}}) \rightarrow 0,
\end{equation}
reflecting the fact that the global state becomes maximally coherent and the two
sectors are perfectly correlated. This CPT-symmetric fixed point determines the
matching conditions for tensor perturbations.

Because the metric is continuous and differentiable across the bounce, tensor modes
must satisfy
\begin{equation}
    h_{ij}(t_{\mathrm{b}}^{-}) = h_{ij}(t_{\mathrm{b}}^{+}),
    \qquad
    \dot{h}_{ij}(t_{\mathrm{b}}^{-}) = \dot{h}_{ij}(t_{\mathrm{b}}^{+}),
\end{equation}
where $t_{\mathrm{b}}^{-}$ and $t_{\mathrm{b}}^{+}$ denote the moments immediately before and after
the bounce. These conditions ensure that tensor perturbations evolve smoothly through
the transition.

The global CPT symmetry imposes an additional constraint. The global state satisfies
\begin{equation}
    \mathrm{CPT}\,\ket{\Psi} = \ket{\Psi},
\end{equation}
which implies for tensor modes
\begin{equation}
    h^{(+)}_{ij}(t,\mathbf{x})
    = h^{(-)}_{ij}(-t,\mathbf{x})\, e^{i\varphi_{\mathrm{CPT}}},
\end{equation}
where $\varphi_{\mathrm{CPT}}$ is a phase determined by the CPT eigenvalue of the global
state. This relation expresses that the visible and mirror tensor modes are time-reversed
images of each other at the bounce.

Away from the bounce, the relaxation parameter grows according to
\begin{equation}
    \dot{\alpha} = \alpha H,
\end{equation}
and approaches its infrared fixed point $\alpha_{\mathrm{IR}}$. This induces a small but
cumulative logarithmic phase difference between the two tensor branches. Writing the
Fourier modes as
\begin{equation}
    h^{(\pm)}(k,t)
    = A_{\pm}(k)\, e^{-i\omega(k)t},
\end{equation}
the CPT relation implies
\begin{equation}
    h^{(+)}(k,t)
    = h^{(-)}(k,-t)\, e^{i\varphi_{\mathrm{CPT}}}.
\end{equation}

The relaxation dynamics modify this relation by introducing a logarithmic phase shift,
\begin{equation}
    \Delta\phi(k)
    = \alpha_{\mathrm{IR}} \ln\!\left( \frac{k}{k_{\mathrm{b}}} \right),
\end{equation}
where $k_{\mathrm{b}}$ is the characteristic bounce scale. This phase shift is the structural
origin of the log-periodic modulation of the tensor spectrum.

Combining the bounce matching conditions with the CPT relation yields the general
form of the tensor mode in the expanding phase,
\begin{equation}
    h(k)
    = h_{0}(k)\left[
        1 + \varepsilon(k)\,
        \cos\!\left(
            \alpha_{\mathrm{IR}} \ln\!\frac{k}{k_{\mathrm{b}}}
            + \varphi_{\mathrm{CPT}}
        \right)
    \right],
\end{equation}
where $\varepsilon(k)$ is a slowly varying amplitude determined by the bounce matching.
This expression already contains the essential structure of the RDCC tensor signature:
a discrete scale invariance encoded in oscillations in $\ln k$.

In the next section we derive this modulation more explicitly by solving the tensor
equation in Fourier space and matching the contracting and expanding solutions across
the bounce.


\begin{figure}[H]
    \centering
    \includegraphics[width=0.78\textwidth]{tensor_interference.pdf}
    \caption{
        Schematic illustration of tensor-sector interference in RDCC.
        The visible and mirror tensor modes are related by the CPT condition
        $h^{(+)}(t,k) = h^{(-)}(-t,k)\, e^{i\varphi_{\mathrm{CPT}}}$ at the bounce.
        Away from the bounce, the relaxation dynamics induce a logarithmic
        phase shift $\Delta\phi(k) = \alpha_{\mathrm{IR}} \ln(k/k_{\mathrm{b}})$.
        Their superposition produces the characteristic interference pattern
        that leads to log-periodic modulation of the tensor spectrum.
    }
    \label{fig:tensor_interference}
\end{figure}



\section{Derivation of Log-Periodic Modulation}

The log-periodic modulation of the stochastic gravitational-wave background (SGWB)
is a structural consequence of the CPT-symmetric two-sector architecture of RDCC.
The key ingredients are: (i) the global tensor equation, (ii) the bounce matching
conditions, and (iii) the CPT relation between the visible and mirror tensor modes.
Together, these imply that tensor modes acquire a logarithmic phase difference whose
magnitude is controlled by the infrared relaxation parameter $\alpha_{\mathrm{IR}}$.

\subsection{Mode Equation in Fourier Space}

We write the tensor perturbation in Fourier space as
\begin{equation}
    h_{ij}(t,\mathbf{x})
    = \sum_{\lambda=+,\times}
      \int \frac{d^{3}k}{(2\pi)^{3}}
      \, h_{\lambda}(t,k)\,
      e^{i\mathbf{k}\cdot\mathbf{x}}\,
      e^{(\lambda)}_{ij}(\mathbf{k}),
\end{equation}
where $e^{(\lambda)}_{ij}$ are the polarization tensors. The mode functions satisfy
\begin{equation}
    \ddot{h}_{\lambda}(k,t)
    + 3H \dot{h}_{\lambda}(k,t)
    + \frac{k^{2}}{a^{2}} h_{\lambda}(k,t)
    = 0,
\end{equation}
in the absence of anisotropic stress. Near the bounce, the Hubble parameter vanishes,
\begin{equation}
    H(t_{\mathrm{b}}) = 0,
\end{equation}
and the relaxation parameter approaches zero,
\begin{equation}
    \alpha(t_{\mathrm{b}}) \rightarrow 0,
\end{equation}
so the two sectors become maximally correlated.

\subsection{Contracting and Expanding Solutions}

In the contracting phase ($t < t_{\mathrm{b}}$), the general solution may be written as
\begin{equation}
    h_{\mathrm{con}}(k,t)
    = A(k)\, e^{-i\omega \ln(k/k_{\mathrm{b}})}
    + B(k)\, e^{+i\omega \ln(k/k_{\mathrm{b}})},
\end{equation}
where $\omega$ is a dimensionless frequency determined by the relaxation dynamics.
The logarithmic dependence arises because the effective mass term scales as
$k^{2}/a^{2}$ and the scale factor evolves approximately exponentially in conformal
time near the bounce.

In the expanding phase ($t > t_{\mathrm{b}}$), the solution takes the form
\begin{equation}
    h_{\mathrm{exp}}(k,t)
    = C(k)\, e^{-i\omega \ln(k/k_{\mathrm{b}})}
    + D(k)\, e^{+i\omega \ln(k/k_{\mathrm{b}})}.
\end{equation}

\subsection{Bounce Matching}

The bounce matching conditions,
\begin{equation}
    h_{\mathrm{con}}(t_{\mathrm{b}}) = h_{\mathrm{exp}}(t_{\mathrm{b}}),
    \qquad
    \dot{h}_{\mathrm{con}}(t_{\mathrm{b}}) = \dot{h}_{\mathrm{exp}}(t_{\mathrm{b}}),
\end{equation}
relate the coefficients $(A,B)$ to $(C,D)$. These relations encode the effect of the
bounce geometry and the vanishing of $\alpha$ at $t_{\mathrm{b}}$.

\subsection{CPT-Induced Phase Structure}

The global CPT symmetry implies
\begin{equation}
    h^{(+)}(k,t)
    = h^{(-)}(k,-t)\, e^{i\varphi_{\mathrm{CPT}}},
\end{equation}
where $\varphi_{\mathrm{CPT}}$ is a constant phase. Away from the bounce, the relaxation
parameter grows toward its infrared fixed point $\alpha_{\mathrm{IR}}$, inducing a logarithmic
phase shift between the two tensor branches,
\begin{equation}
    \Delta\phi(k)
    = \alpha_{\mathrm{IR}} \ln\!\left( \frac{k}{k_{\mathrm{b}}} \right).
\end{equation}

This phase shift is the structural origin of the log-periodic modulation. The visible
sector perceives the superposition
\begin{equation}
    h(k)
    = h^{(+)}(k) + h^{(-)}(k),
\end{equation}
which yields
\begin{equation}
    h(k)
    = h_{0}(k)\left[
        1 + \varepsilon(k)\,
        \cos\!\left(
            \alpha_{\mathrm{IR}} \ln\!\frac{k}{k_{\mathrm{b}}}
            + \varphi_{\mathrm{CPT}}
        \right)
    \right],
\end{equation}
where $\varepsilon(k)$ is a slowly varying amplitude determined by the bounce matching.

\subsection{Discrete Scale Invariance}

The modulation term is periodic in $\ln k$ with period
\begin{equation}
    \Delta(\ln k)
    = \frac{2\pi}{\alpha_{\mathrm{IR}}},
\end{equation}
reflecting a discrete scale invariance of the tensor spectrum. This is the hallmark of
RDCC: the modulation frequency is fixed entirely by the infrared relaxation parameter,
\begin{equation}
    \omega_{\mathrm{mod}} = \alpha_{\mathrm{IR}}.
\end{equation}

The resulting tensor spectrum,
\begin{equation}
    P_{h}(k)
    = P_{0}(k)\left[
        1 + \varepsilon(k)\,
        \cos\!\left(
            \alpha_{\mathrm{IR}} \ln\!\frac{k}{k_{\mathrm{b}}}
            + \varphi_{\mathrm{CPT}}
        \right)
    \right],
\end{equation}
exhibits a blue tilt, a millihertz-scale peak, and a log-periodic modulation whose
frequency is controlled by a single parameter. This structure provides a sharply
falsifiable prediction for upcoming gravitational-wave observatories.

\begin{figure}[H]
    \centering
    \includegraphics[width=0.78\textwidth]{log_periodic_modulation.pdf}
    \caption{
        One-dimensional illustration of the RDCC log-periodic modulation.
        The tensor spectrum takes the form
        $P_{h}(k) = P_{0}(k)\left[1 + \varepsilon(k)\,
        \cos\!\left(\alpha_{\mathrm{IR}} \ln(k/k_{\mathrm{b}})
        + \varphi_{\mathrm{CPT}}\right)\right]$.
        The oscillations in $\ln k$ reflect a discrete scale invariance with
        period $\Delta(\ln k) = 2\pi/\alpha_{\mathrm{IR}}$, a structural
        consequence of the CPT-symmetric two-sector architecture.
    }
    \label{fig:log_periodic_modulation}
\end{figure}


\section{Observational Consequences}

The tensor signature of RDCC is characterized by three structural features:
(i) a blue-tilted spectrum, (ii) a millihertz-scale peak, and (iii) a log-periodic
modulation with frequency proportional to the infrared relaxation parameter
$\alpha_{\mathrm{IR}}$. These features span more than twenty decades in frequency and
are accessible to a broad range of gravitational-wave observatories.

\subsection{Low-Frequency Regime: CMB Experiments}

At the lowest frequencies, the RDCC tensor spectrum contributes to the CMB
$B$-mode polarization. The log-periodic modulation is strongly suppressed in this
regime, but the overall amplitude and tilt remain sensitive to the relaxation
parameter. Experiments such as LiteBIRD and CMB-S4 probe the combination
\begin{equation}
    \Omega_{\mathrm{GW}}(f_{\mathrm{CMB}})
    \propto \alpha_{\mathrm{IR}}\, f_{\mathrm{CMB}}^{\,n_{T}},
\end{equation}
where $n_{T} > 0$ is the blue tensor tilt. A detection of a blue-tilted tensor
spectrum at CMB scales would be incompatible with standard inflationary
predictions and would provide indirect support for RDCC.

\subsection{Millihertz Regime: LISA}

The peak of the RDCC tensor spectrum lies naturally in the millihertz band,
\begin{equation}
    f_{\mathrm{peak}} \sim \mathcal{O}(\mathrm{mHz}),
\end{equation}
making LISA the prime detector for the RDCC signal. The log-periodic modulation
is most pronounced near the peak, where the interference between the visible and
mirror tensor sectors is strongest. The modulation period,
\begin{equation}
    \Delta(\ln f)
    = \frac{2\pi}{\alpha_{\mathrm{IR}}},
\end{equation}
is large for $\alpha_{\mathrm{IR}} \sim 10^{-2}$, implying that only a fraction of a full
oscillation is visible within the LISA band. Even a partial oscillation would
constitute a distinctive signature.

\subsection{Mid-Frequency Regime: DECIGO and BBO}

Future space-based detectors such as DECIGO and BBO offer significantly broader
frequency coverage than LISA. Their sensitivity to the mid-frequency regime allows
them to probe multiple cycles of the log-periodic modulation. This makes them
ideal instruments for testing the discrete scale invariance predicted by RDCC.
A detection of repeated oscillations in $\ln f$ would provide a direct measurement
of $\alpha_{\mathrm{IR}}$ and would serve as a definitive test of the framework.

\subsection{High-Frequency Regime: Ground-Based Detectors}

Ground-based interferometers such as the Einstein Telescope (ET) and Cosmic
Explorer (CE) probe the high-frequency tail of the RDCC spectrum. In this regime,
the blue tilt enhances the amplitude relative to inflationary predictions. Although
the log-periodic modulation is less pronounced at high frequencies, the overall
spectral shape remains a powerful discriminator. The combination of LISA,
DECIGO/BBO, and ET/CE provides a multi-band observational strategy capable of
testing the full RDCC tensor signature across more than twenty decades in
frequency.

\subsection{Summary of Observational Signatures}

The observational consequences of the RDCC tensor sector may be summarized as:
\begin{itemize}
    \item a blue-tilted SGWB with $n_{T} > 0$,
    \item a millihertz-scale peak set by the bounce scale,
    \item a log-periodic modulation with frequency $\omega_{\mathrm{mod}} = \alpha_{\mathrm{IR}}$,
    \item and a discrete scale invariance encoded in oscillations in $\ln f$.
\end{itemize}
These features arise from the same structural mechanism: the CPT-symmetric
two-sector architecture and the relaxation dynamics governed by
$\alpha_{\mathrm{IR}}$. Their simultaneous detection would provide strong evidence
for the RDCC framework, while their absence would falsify it.

\section{Falsifiability}

The tensor sector of RDCC provides a sharply falsifiable prediction. Its one-parameter
structure implies that all deviations from $\Lambda$CDM are correlated and controlled by
the infrared relaxation parameter $\alpha_{\mathrm{IR}}$. The log-periodic modulation of the
stochastic gravitational-wave background (SGWB) is the most distinctive and model-specific
signature of this framework.

\subsection{Absence of Log-Periodic Modulation}

The presence of oscillations in $\ln k$ or $\ln f$ is a direct consequence of the CPT-symmetric
two-sector structure. If future gravitational-wave observations reveal a smooth tensor
spectrum without oscillatory features, RDCC would be ruled out. No alternative mechanism
within the framework can reproduce a featureless spectrum.

\subsection{Inconsistent Amplitude Scaling}

The SGWB amplitude scales as
\begin{equation}
    \Omega_{\mathrm{GW}} \propto \alpha_{\mathrm{IR}}^{\,2},
\end{equation}
while the scalar spectral tilt satisfies
\begin{equation}
    n_{s} - 1 = -2\alpha_{\mathrm{IR}}.
\end{equation}
A measured tensor amplitude inconsistent with the value of $\alpha_{\mathrm{IR}}$ inferred from
the scalar sector would falsify the framework. This correlation is exact within RDCC and
does not depend on model-specific assumptions.

\subsection{Incorrect Peak Frequency}

The peak frequency of the tensor spectrum is determined by the bounce scale and lies
naturally in the millihertz range. A significantly different peak frequency would be
incompatible with the cuscuton-induced bounce dynamics and would rule out the RDCC
interpretation of the SGWB.

\subsection{Violation of Correlated Predictions}

The one-parameter nature of RDCC implies the following consistency relations:
\begin{equation}
    n_{s} - 1 = -2\alpha_{\mathrm{IR}},
    \qquad
    \Delta N_{\mathrm{eff}} \sim \alpha_{\mathrm{IR}},
    \qquad
    \Omega_{\mathrm{GW}} \propto \alpha_{\mathrm{IR}}^{\,2}.
\end{equation}
Any inconsistency among these observables would falsify the framework. Conversely, a
coherent detection across all channels would provide strong evidence for the relaxation-driven
structure of the Universe.

\subsection{Summary}

The tensor sector offers a decisive and experimentally accessible test of RDCC. The detection
of log-periodic modulation in the SGWB, together with the correlated predictions for the
scalar tilt and dark-radiation excess, would constitute compelling evidence for a CPT-symmetric
two-sector Universe. The absence of these signatures would rule out the framework.

\section{Conclusion}

The tensor sector of Relaxation-Driven Cyclic Cosmology (RDCC) provides a uniquely
sharp and experimentally accessible window into the CPT-symmetric structure of the
Universe. Unlike inflationary scenarios, where tensor modes originate from stochastic
amplification at horizon exit, RDCC generates tensor perturbations through CPT-consistent
matching across a nonsingular bounce. The bipartite Hilbert-space structure implies that
visible and mirror tensor sectors interfere, producing a log-periodic modulation of the
stochastic gravitational-wave background (SGWB).

The resulting tensor spectrum exhibits three structural features:
\begin{itemize}
    \item a blue tilt originating from the cuscuton-induced bounce,
    \item a millihertz-scale peak set by the bounce scale,
    \item and a log-periodic modulation with frequency
    $\omega_{\mathrm{mod}} = \alpha_{\mathrm{IR}}$.
\end{itemize}

Each of these features arises directly from the one-parameter nature of RDCC and the
global CPT symmetry of the cosmological state. The modulation frequency is fixed by the
infrared relaxation parameter, reflecting the fact that the same mechanism governs the
scalar tilt, the dark-radiation excess, and the tensor amplitude.

The observational prospects are exceptionally strong. LISA probes the peak region,
DECIGO and BBO can resolve multiple oscillation cycles, and ground-based detectors test
the high-frequency tail. CMB experiments constrain the low-frequency limit. Together,
these instruments provide a multi-band strategy capable of confirming or falsifying the
RDCC tensor signature.

The tensor sector therefore stands as the decisive experimental test of RDCC. A detection
of log-periodic modulation would constitute compelling evidence for a CPT-symmetric
two-sector Universe, while its absence would rule out the framework.

\begin{thebibliography}{9}

\bibitem{RDCC-Framework}
M.~Lehmann,
\textit{Relaxation-Driven Cyclic Cosmology (RDCC) Framework},
Independent Research (2026).

\bibitem{Companion-I}
M.~Lehmann,
\textit{Companion I: Gravity as the Global Bridge in RDCC},
Independent Research (2026).

\bibitem{Cuscuton}
N.~Afshordi, D.~J.~H.~Chung, G.~Geshnizjani,
\textit{Cuscuton: A Causal Field Theory with an Infinite Speed of Sound},
Phys.\ Rev.\ D 75, 083513 (2007).

\bibitem{LISA}
P.~Amaro-Seoane et al.,
\textit{Laser Interferometer Space Antenna},
arXiv:1702.00786.

\bibitem{DECIGO}
S.~Kawamura et al.,
\textit{The Japanese Space Gravitational Wave Antenna: DECIGO},
Class.\ Quant.\ Grav.\ 28, 094011 (2011).

\bibitem{ET}
M.~Punturo et al.,
\textit{The Einstein Telescope: A Third-Generation Gravitational Wave Observatory},
Class.\ Quant.\ Grav.\ 27, 194002 (2010).

\end{thebibliography}

\end{document}


\end{document}
